\section{Discussion}
\label{sec:discussion}

\subsection{What the Results Establish}

The framework's contribution is narrower than the one claimed for it at the outset, and
better supported. Three findings survive the evaluation of Section~4.

The first concerns the separation of the allocation decision. Splitting the problem into
relative composition and total risk, rather than allowing a single constraint set to
determine both, makes each part independently testable. It is this separation that
allows the ablation in Table~\ref{tab:ablation} to be constructed at all: without it, removing the regime
layer would change both what is held and how much of it is held, and the resulting
comparison would confound the two.

The second is that regime conditioning earns its place through timing, not through
level. At identical average exposure, allocating risk according to the identified state
rather than uniformly improves the Sharpe ratio from 0.795 to 0.878 and maximum
drawdown from $-22.26$ to $-18.98$ per cent, with no significant difference in mean
return ($p = 0.695$). This is the cleanest evidence in the paper, and it is available only
because the control was constructed at matched risk. The corresponding cost is also
measured: against a fully invested variant the framework forgoes approximately 1.3
percentage points of annual return, significantly so, in exchange for 3.6 percentage
points of maximum drawdown.

The third is that the form of the risk budget matters more than its calibration. The
four relative budgets in Table~\ref{tab:budget} produce Sharpe ratios within 0.020 of one another,
while the absolute-target formulation---identical in intent---is inert, binding in one
rebalance out of 224. An absolute volatility target is not scale-free, and a
specification transferred between universes without re-calibration can silently disable
the mechanism it was written to implement. This is a practical warning that generalises
beyond the present framework.

\subsection{What the Results Do Not Establish}

Two claims that the framework was designed to support are not supported.

The signal fusion layer does not contribute. Its removal improves the Sharpe ratio from
0.878 to 0.908 and leaves maximum drawdown unchanged, and neither movement is
significant. The information coefficients in Table~\ref{tab:ic} explain why: both composites are
significantly negatively related to next-session cross-sectional returns over the full
sample. The regime-conditional weighting schedule has no measured basis either. It was
specified on the hypothesis that market-state information becomes relatively more
informative under stress; the point estimates run the other way and the calm-to-stress
contrast is not significant ($t = -1.46$, $p = 0.145$), so the schedule is weighting on a
difference the data do not establish. Splitting the sample at the holdout boundary
sharpens the reading: the fused coefficient is $-0.039$ before 2024 and $+0.009$ after,
so even the negative relationship is era-specific and could not have been exploited by
inverting the composites.

The improvement in risk-adjusted return is not statistically significant. The Sharpe
difference of $+0.250$ carries a bootstrap interval of $[-0.102, 0.574]$. Eighteen years
of daily returns contain only a handful of independent stress events, and a claim about
behaviour during market stress cannot be established at conventional significance from
four episodes. Adding the 665 sessions of the holdout narrowed the interval without
moving it clear of zero. The consistency of the result across 18 of 18 specifications,
18 of 19 calendar years and 4 of 4 episodes is the stronger evidence, and it is evidence of
robustness rather than of significance. The two should not be conflated, and this paper
does not conflate them.

\subsection{Implications for Practice}

For a practitioner the results suggest a specific and limited use. A regime-conditional
risk budget is a reasonable component of a multi-asset allocation process where the
objective includes drawdown control and where some return may be sacrificed for it. It
should be implemented in relative rather than absolute form, evaluated against a
volatility-matched benchmark rather than an unlevered one, and not expected to help in
a universe without genuine dispersion in asset volatility.

The results also caution against a common design pattern. The framework as originally
specified contained two layers whose contribution had never been isolated, one of which
was inert by construction and the other of which was pointed in the wrong direction.
Neither defect was visible in aggregate performance, and both were found only by
ablation at matched risk. Architectures assembled from plausible components should be
disassembled before they are trusted.

Finally, the turnover result in Table~\ref{tab:turnover} has a practical reading. The assumed
transaction cost enters the objective as a penalty as well as the return as a
deduction, so setting it to zero does not produce the best achievable result---it
produces a noisier weight vector and a lower Sharpe ratio. A modest cost assumption acts
as a regulariser, and there is no reason to model a frictionless world even when
evaluating one.

\subsection{Limitations}

Several limitations bound what can be concluded.

The universe is small. Ten exchange-traded funds are enough to demonstrate the
mechanism but not to establish that it generalises to larger or differently composed
universes, and the position cap of 25 per cent binds on average, which means the
optimiser's freedom is partly determined by the constraint rather than by the risk
model.

The evaluation is a single historical path. Eighteen years is a long sample by the
standards of this literature and a short one for inference about tail events, and the
four stress episodes are not independent draws from a stationary process.

The signals are market-derived proxies rather than textual sentiment. No news, filings
or social media text is used at any point, as stated in Section~3.3. The null reported
in Section~4.7 is therefore a null about these proxies, and it is not evidence that
textual sentiment would fail in the same role. Establishing that would require a
time-stamped news corpus with coverage of the traded instruments, which for broad
exchange-traded funds does not exist at the instrument level.

The regime model is deliberately simple. A three-component Gaussian mixture on two
features is interpretable and stable, but it imposes conditional normality within
states, treats observations as conditionally independent given the state, and---as
Table~\ref{tab:regimefit} shows---is not the specification either information criterion would select.

The risk-free rate is taken as zero throughout, and cash is assumed to earn nothing.
Over a sample containing both the near-zero rates of 2009--2015 and the higher rates of
2023, this understates the return on the framework's cash holding and therefore
understates its performance relative to the fully invested benchmark. The direction of
this bias is against the framework, which is the conservative choice, but it is a
simplification.

\subsection{Future Work}
\label{sec:future}

Four extensions follow directly. The first is textual sentiment, applied to a universe
where an instrument-level news corpus exists, which would convert the proxy null of
Section~4.7 into a genuine test of the fusion hypothesis. The second is a larger and
more heterogeneous universe, including commodities, currencies and non-US fixed income,
which would test whether the volatility-dispersion condition identified in
Section~4.11 is the binding requirement it appears to be. The third is learning the
risk-budget multipliers rather than fixing them, for which reinforcement learning is a
natural candidate given the sequential structure of the problem, subject to the
overfitting concerns documented by \citet{harvey15} and \citet{valls26}. The fourth is
replacing the mixture model with a formulation that represents state persistence
explicitly, such as a hidden Markov or jump model \citep{li25}, which would remove the
conditional-independence assumption noted above.
